% =====================================================================
%  Seance 6 -- D. TP : du modele au service + evaluation blanche (23--28)
% =====================================================================
\section{TP « De la forêt au service » et évaluation blanche}

\begin{frame}{Le TP : objectif et déroulé (2 h)}
\begin{columns}[t]
\begin{column}{0.52\textwidth}
\begin{kbox}[Ce que vous produisez à 18 h]
\begin{enumerate}
\item \textbf{B1} --- \file{nosys\_pipeline.joblib} + un script de scoring
qui charge le pipeline et score un nouveau fichier de métriques (30 min).
\item \textbf{B2} --- un rapport de latence : moyenne et p99 sur
\textbf{1 000 prédictions réelles}, et le raisonnement de capacité (30 min).
\item \textbf{B3} --- une fiche MLOps d'une page : plan de retraining,
métriques surveillées, procédure de rollback (30 min).
\item \textbf{B4} --- évaluation blanche : 20 questions couvrant les 6
séances, corrigé fourni (30 min).
\end{enumerate}
\end{kbox}
\end{column}
\begin{column}{0.44\textwidth}
\begin{ebox}[Le pipeline officiel du cours]
Tout part de \file{nosys\_features.csv} et du split temporel strict :
train $<$ 2026-09-22 (24 192 fenêtres, 56 positives), test $\ge$
(10 368, 64 positives). \textbf{Aucune ligne du futur dans le train.}
\end{ebox}
\begin{wbox}[Le test de vraie fin]
Si votre B1 sauvegarde \textbf{la standardisation en même temps que la
forêt}, c'est gagné : un pipeline partiel (modèle seul) produit des scores
faux sur un nouveau fichier. C'est le piège classique de l'industrialisation.
\end{wbox}
\end{column}
\end{columns}
\end{frame}

\begin{frame}[fragile]{TP B1 --- sauvegarder le pipeline avec joblib (30 min)}
\small
\begin{codebox}[B1. Le pipeline complet, pas seulement la forêt]
\begin{lstlisting}
from sklearn.ensemble import RandomForestClassifier
from sklearn.preprocessing import StandardScaler
from sklearn.pipeline import make_pipeline
import joblib

pipe = make_pipeline(StandardScaler(),
        RandomForestClassifier(n_estimators=200, min_samples_leaf=2,
            class_weight="balanced", random_state=5534))
pipe.fit(Xtr, ytr)                 # train < 2026-09-22
joblib.dump(pipe, "out/nosys_pipeline.joblib")

pipe = joblib.load("out/nosys_pipeline.joblib")   # dans le service
nouveau = pd.read_csv("data/nouveaux_metriques.csv")
scores = pipe.predict_proba(nouveau[variables])[:, 1]
\end{lstlisting}
\end{codebox}
\small\begin{abox}[Livrable B1]
\file{nosys\_pipeline.joblib} (dans \file{out/}, jamais dans \file{data/})
+ un script \py{scorer.py} qui prend un CSV de métriques et écrit les
scores. Test : scorer le train doit redonner les probabilités d'avant
sauvegarde.
\end{abox}
\end{frame}

\begin{frame}{TP B2 --- mesurer la latence réelle de scoring (30 min)}
\small
\begin{center}
\includegraphics[width=0.62\textwidth]{f34_latence.png}
\end{center}
\begin{columns}[t]
\begin{column}{0.5\textwidth}
\begin{ebox}[Résultats réels du script du cours (1 000 prédictions)]
Moyenne : \textbf{59 ms} par fenêtre --- p99 : \textbf{105 ms}. Par lots de
100 : \textbf{94 279 fenêtres/min} par cœur. Les chiffres de votre machine
seront vos propres mesures --- c'est ça le livrable.
\end{ebox}
\end{column}
\begin{column}{0.46\textwidth}
\begin{qbox}[Le raisonnement de capacité à rédiger]
NOSYS produit \textbf{1 152 fenêtres/jour} (8 serveurs). Un cœur à 1 000
fenêtres/min en individuel : quelle marge ? Et pourquoi le scoring par lots
est-il $\approx$ 90 fois plus rapide par fenêtre ? (indice : overhead
d'appel, parallélisme interne des arbres)
\end{qbox}
\end{column}
\end{columns}
\end{frame}

\begin{frame}{TP B3 --- la fiche MLOps d'une page (30 min)}
\small
\begin{columns}[t]
\begin{column}{0.52\textwidth}
\small\begin{abox}[Les trois rubriques imposées]
\begin{enumerate}
\item \textbf{Plan de retraining} : fréquence de base (ex. hebdomadaire) et
\textbf{déclencheurs chiffrés} --- ex. « proba moyenne $>$ 0,005 pendant
48 h », « taux d'alertes $>$ 0,2 \%», « AUC $<$ 0,70 dès que mesurable ».
\item \textbf{Métriques surveillées} : AUC, taux d'alertes, latence moyenne
+ p99, distribution des entrées --- avec les valeurs réelles du B2 et de f29.
\item \textbf{Procédure de rollback} : garder les N dernières versions ; si
le taux d'alertes double après bascule, revenir à la version précédente en
moins de 10 minutes.
\end{enumerate}
\end{abox}
\end{column}
\begin{column}{0.44\textwidth}
\small\begin{kbox}[Barème B3 (sur 7 pts de la note TP)]
\begin{itemize}
\item Déclencheurs \textbf{chiffrés} et justifiés par les mesures --- 3 pts
\item Métriques complètes avec valeurs réelles --- 2 pts
\item Rollback testable (qui, quoi, délai) --- 2 pts
\end{itemize}
\end{kbox}
\begin{wbox}[Le critère de réussite]
Une fiche MLOps est bonne quand un collègue qui n'a pas suivi le cours
peut faire tourner et surveiller le modèle \textbf{avec elle seule}.
\end{wbox}
\end{column}
\end{columns}
\end{frame}

\begin{frame}{TP B4 --- évaluation blanche : 20 questions (30 min)}
\small
\begin{columns}[t]
\begin{column}{0.52\textwidth}
\begin{ebox}[Format]
QCM identique à l'examen final : \textbf{20 questions}, une par grande
notion : 3 sur les données (séances 1--2), 4 sur classification (séance 3),
3 sur les séries temporelles (séance 4), 3 sur les logs (séance 5),
7 sur MLOps et gouvernance (séance 6).
\end{ebox}
\end{column}
\begin{column}{0.44\textwidth}
\small\begin{kbox}[Exemples de questions]
\begin{enumerate}
\item Le split strict sépare train et test à quelle date ? \hfill (22/09)
\item Pourquoi p99 et pas seulement la moyenne de latence ?\\
(une moyenne cache les rares prédictions très lentes)
\item Que surveille-t-on sur une tranche sans incident ?\\
(la distribution des scores, pas l'AUC)
\end{enumerate}
\end{kbox}
\begin{solbox}[Corrigé]
Distribué après les 30 min --- chaque réponse renvoie à la séance et à la
frame concernée.
\end{solbox}
\end{column}
\end{columns}
\end{frame}
